% === BEGIN SEGMENT 1 ===
\documentclass[journal,twoside]{IEEEtran}

\usepackage{amsmath}
\usepackage{amssymb}
\usepackage[ruled,vlined]{algorithm2e}
\usepackage{booktabs}
\usepackage{graphicx}
\usepackage{hyperref}
\usepackage{cite}
\usepackage{tikz}
\usepackage{pgfplots}
\pgfplotsset{compat=1.18}
\usepackage{multirow}
\usepackage{xcolor}
\usepackage{balance}

\hypersetup{colorlinks=true,linkcolor=blue,citecolor=blue,urlcolor=blue}

\begin{document}

\title{RobustIDPS: An Adversarially Robust Intrusion Detection System with Continual Learning and Constrained Reinforcement Learning for Autonomous Response}

\author{
  \IEEEauthorblockN{Roger Nick Anaedevha}
  \IEEEauthorblockA{National Research Nuclear University MEPhI\\
  Moscow, Russia\\
  roger@robustidps.ai}
}

\maketitle

% ---------------------------------------------------------------
% ABSTRACT
% ---------------------------------------------------------------
\begin{abstract}
Modern network intrusion detection and prevention systems (IDPS) face compounding
challenges: rapidly evolving attack landscapes, adversarial evasion of
machine-learning classifiers, distributional drift in production traffic, and the
need for timely autonomous response without excessive false-positive disruption.
We present \textsc{RobustIDPS}, a unified framework that addresses all four
challenges simultaneously. At its core, a seven-branch surrogate ensemble
augmented with Monte Carlo dropout uncertainty estimation ($T{=}20$ stochastic
forward passes) classifies 83 engineered features into 34 attack classes.
Continual learning is achieved through Elastic Weight Consolidation (EWC)
combined with experience replay, sharing a unified Fisher information framework
with adversarial robustness via a balancing coefficient $\beta{=}0.7$.
A Constrained Policy Optimization (CPO) agent maps detector outputs to five
graduated response actions---Monitor, Rate-Limit, Reset, Block, and
Quarantine---while satisfying user-specified safety constraints on
false-positive-induced service disruption. Adversarial robustness is
systematically evaluated against six attack families: FGSM, PGD, C\&W,
DeepFool, Gaussian noise, and label-masking poisoning. For distributed
deployment, we introduce FedGTD, a Byzantine-resilient federated learning
protocol based on geometric trimmed deviation that tolerates up to 30\%
malicious participants. Temporal traffic dynamics are captured by a
Stochastic Differential Equation--Temporal Graph Neural Network (SDE-TGNN).
Comprehensive experiments on CICIDS2017, NSL-KDD, CIC-IoT-2023, and
UNSW-NB15 demonstrate that \textsc{RobustIDPS} achieves a macro-averaged
F1 score of 0.967 under clean conditions, degrades by less than 3.1\% under
the strongest PGD-40 attack ($\epsilon{=}0.03$), and maintains above 0.94 F1
after five sequential concept-drift episodes without catastrophic forgetting.
\end{abstract}

\begin{IEEEkeywords}
Intrusion detection, adversarial robustness, continual learning, constrained
reinforcement learning, federated learning, graph neural networks
\end{IEEEkeywords}

% ---------------------------------------------------------------
% I. INTRODUCTION
% ---------------------------------------------------------------
\section{Introduction}
\label{sec:introduction}

The proliferation of sophisticated cyber threats has rendered signature-based
intrusion detection systems (IDS) increasingly inadequate for protecting
enterprise and critical-infrastructure networks~\cite{khraisat2019survey}.
Machine learning (ML) approaches have demonstrated considerable promise in
detecting both known and zero-day attacks by learning discriminative
representations from network telemetry~\cite{buczak2016survey,
aldweesh2020deep}. Yet, the deployment of ML-based network intrusion detection
systems (NIDS) in adversarial environments exposes fundamental vulnerabilities
that remain largely unresolved.

First, ML classifiers are susceptible to adversarial examples---small,
deliberately crafted perturbations that cause misclassification while remaining
within the feasible feature space of network
traffic~\cite{goodfellow2015explaining, carlini2017towards}. An attacker who
understands the feature extraction pipeline can evade detection with
perturbations imperceptible to rule-based audit~\cite{hashemi2019towards}.
Second, network traffic distributions are inherently non-stationary: new
applications, protocols, and attack toolkits continuously shift the data
manifold, causing concept drift that degrades deployed
models~\cite{jordaney2017transcend}. Periodic retraining from scratch is
expensive and discards previously learned structure. Third, even an accurate
detector is of limited operational value without a principled response
mechanism; manual incident response cannot scale to the volume and velocity
of modern attacks~\cite{nguyen2021deep}. Fourth, privacy regulations and
organizational boundaries demand that detection models be trainable across
distributed sites without centralizing raw traffic, yet na\"ive federated
learning is vulnerable to Byzantine participants submitting poisoned
gradients~\cite{blanchard2017machine}.

Despite a rich literature on individual aspects---adversarial ML for
security~\cite{apruzzese2023real}, continual
learning~\cite{delange2021continual}, reinforcement learning (RL) for cyber
defense~\cite{nguyen2021deep}, and federated
learning~\cite{mcmahan2017communication}---no existing system integrates all
four capabilities within a coherent, theoretically grounded framework. This
fragmentation forces practitioners to assemble ad-hoc pipelines whose
component interactions are neither analyzed nor optimized.

\textsc{RobustIDPS} fills this gap. Our key insight is that the Fisher
information matrix, central to both EWC-based continual
learning~\cite{kirkpatrick2017overcoming} and curvature-aware adversarial
training~\cite{wu2020adversarial}, can serve as a \emph{unified} regularization
backbone. By sharing Fisher estimates across the continual-learning and
adversarial-robustness objectives through a single balancing coefficient
$\beta$, we eliminate redundant computation and expose a principled
accuracy--robustness--plasticity trade-off.

The principal contributions of this work are:

\begin{enumerate}
  \item A \textbf{seven-branch surrogate ensemble} architecture with MC dropout
        uncertainty quantification ($T{=}20$), providing calibrated confidence
        estimates that inform both the RL response agent and the federated
        aggregation protocol.

  \item A \textbf{unified Fisher information framework} ($\beta{=}0.7$) that
        jointly governs EWC continual learning with experience replay and
        curvature-aware adversarial training, enabling graceful
        accuracy--robustness--plasticity trade-offs without catastrophic
        forgetting.

  \item A \textbf{Constrained Policy Optimization (CPO)} response agent that
        maps detector outputs to five graduated actions (Monitor, Rate-Limit,
        Reset, Block, Quarantine) while satisfying hard safety constraints on
        service disruption from false positives.

  \item Systematic \textbf{adversarial robustness} evaluation against six
        attack families (FGSM, PGD, C\&W, DeepFool, Gaussian noise,
        label-masking poisoning), demonstrating less than 3.1\% F1 degradation
        under PGD-40 at $\epsilon{=}0.03$.

  \item \textbf{FedGTD}, a Byzantine-resilient federated aggregation protocol
        based on geometric trimmed deviation that tolerates up to 30\% malicious
        participants while preserving convergence guarantees.

  \item An \textbf{SDE-TGNN} module that captures temporal traffic dynamics
        through stochastic differential equations over a temporal graph,
        improving detection of slow-probe and multi-stage attacks.
\end{enumerate}

We evaluate \textsc{RobustIDPS} on four widely used benchmarks---CICIDS2017,
NSL-KDD, CIC-IoT-2023, and UNSW-NB15---spanning 83 engineered features and 34
attack classes, and demonstrate state-of-the-art results across clean accuracy,
adversarial robustness, continual-learning stability, and response optimality.

% ---------------------------------------------------------------
% II. RELATED WORK
% ---------------------------------------------------------------
\section{Related Work}
\label{sec:related_work}

\subsection{ML-Based Intrusion Detection Systems}
\label{sec:rw_ml_ids}

Classical ML approaches to network intrusion detection include random
forests~\cite{zhang2008random}, support vector
machines~\cite{mukkamala2002intrusion}, and gradient-boosted
trees~\cite{dhaliwal2018effective}. While effective on stationary benchmarks,
these methods rely on hand-crafted features and exhibit limited capacity to
model complex inter-feature dependencies. Deep learning architectures---CNNs
applied to flow images~\cite{wang2017malware}, LSTMs over packet
sequences~\cite{yin2017deep}, and autoencoders for anomaly
detection~\cite{mirsky2018kitsune}---have progressively improved detection
performance. More recently, transformer-based models have shown strong results
on CICIDS2017 and UNSW-NB15~\cite{wu2022rtids, ferriyan2022deep}. Ensemble
strategies that combine heterogeneous base learners further improve robustness
to feature noise~\cite{gao2019adaptive}. However, none of these works
systematically address adversarial evasion, concept drift, or autonomous
response selection within a single framework.

Graph neural networks (GNNs) have gained attention for host-level and
network-level intrusion detection by representing communication patterns as
graphs~\cite{lo2022graphsage, chang2021graph}. Temporal extensions such as
EvolveGCN~\cite{pareja2020evolvegcn} and TGAT~\cite{xu2020inductive} model
time-varying topologies. Our SDE-TGNN advances this line by incorporating
stochastic differential equations to capture the continuous-time evolution of
traffic graphs under uncertainty.

\subsection{Continual Learning for Non-Stationary Environments}
\label{sec:rw_cl}

Continual learning aims to acquire new knowledge without catastrophic
forgetting of previous tasks~\cite{delange2021continual}. Three dominant
strategies exist: regularization-based methods such as EWC~\cite{kirkpatrick2017overcoming}
and Synaptic Intelligence~\cite{zenke2017continual}, replay-based methods
including experience replay~\cite{rolnick2019experience} and generative
replay~\cite{shin2017continual}, and architecture-based methods like
Progressive Neural Networks~\cite{rusu2016progressive}. In the cybersecurity
domain, Andresini et al.~\cite{andresini2021insomnia} applied class-incremental
learning to malware detection, and Barbero et al.~\cite{barbero2022transcending}
demonstrated that EWC mitigates drift-induced degradation in NIDS. However,
existing works treat continual learning in isolation from adversarial robustness.
\textsc{RobustIDPS} unifies these objectives by sharing Fisher information
between the EWC penalty and curvature-aware adversarial training, eliminating
the need for independent regularization hyperparameters.

\subsection{RL for Autonomous Cyber Defense}
\label{sec:rw_rl}

Reinforcement learning has been explored for automated network defense in
simulated environments such as CyberBattleSim~\cite{msft2021cyberbattlesim}
and CAGE Challenge~\cite{cage2022}. Nguyen et al.~\cite{nguyen2021deep}
surveyed deep RL applications in cybersecurity, noting the gap between
simulated success and real-world deployment. Existing approaches typically use
unconstrained policy gradient methods~\cite{elderman2017adversarial,
foley2023autonomous}, which may select aggressive responses (e.g., full
quarantine) that cause unacceptable service disruption. Constrained Markov
Decision Processes (CMDPs)~\cite{altman1999constrained} provide a principled
framework for bounding such costs. Achiam et al.~\cite{achiam2017constrained}
proposed CPO, which guarantees near-constraint satisfaction at every policy
update. We adapt CPO to the IDPS response domain, defining cost functions over
false-positive-induced disruption and demonstrating that constraint enforcement
yields superior operational trade-offs compared to unconstrained baselines
and reward-shaping heuristics.

\subsection{Adversarial Robustness in NIDS}
\label{sec:rw_adv}

Adversarial attacks on NIDS have been demonstrated across both feature-space
and problem-space threat models~\cite{hashemi2019towards, apruzzese2023real}.
Goodfellow et al.~\cite{goodfellow2015explaining} introduced FGSM, and
Madry et al.~\cite{madry2018towards} proposed PGD as a stronger iterative
variant. Carlini and Wagner~\cite{carlini2017towards} developed optimization-based
attacks that minimize perturbation magnitude. In the IDS context,
Peng et al.~\cite{peng2019adversarial} showed that adversarial training
improves robustness on NSL-KDD, while Debicha et al.~\cite{debicha2023adv}
evaluated GAN-based evasion on CICIDS2017. A persistent limitation is that
adversarial training degrades clean accuracy~\cite{tsipras2019robustness}
and interacts poorly with continual-learning regularization. Our unified Fisher
framework addresses this by modulating the curvature penalty jointly across
both objectives, achieving Pareto-superior accuracy--robustness trade-offs
(Section~\ref{sec:fisher_framework}).

Label-flipping and data poisoning attacks represent a complementary threat
axis~\cite{biggio2012poisoning, shafahi2018poison}. We evaluate
\textsc{RobustIDPS} against label-masking poisoning, in which an adversary
corrupts a fraction of training labels during federated rounds, and show that
FedGTD's geometric trimming neutralizes up to 30\% poisoned participants
without significant accuracy loss.

\subsection{Federated Learning for Distributed IDS}
\label{sec:rw_fl}

Federated learning enables collaborative model training without sharing raw
data~\cite{mcmahan2017communication}. In the IDS domain, Nguyen et al.~\cite{nguyen2019diot}
applied federated learning to IoT anomaly detection, and
Popoola et al.~\cite{popoola2021federated} evaluated FedAvg on CICIDS2017
partitions. Byzantine robustness remains a critical concern: Blanchard
et al.~\cite{blanchard2017machine} proposed Krum, while Yin et al.~\cite{yin2018byzantine}
introduced coordinate-wise trimmed mean. More recent protocols include
Bulyan~\cite{guerraoui2018hidden} and centered clipping~\cite{karimireddy2022byzantinerobust}.

Existing Byzantine-resilient aggregation rules operate independently on
each parameter coordinate, which can be exploited by coordinated attacks that
distribute malicious gradients across dimensions~\cite{baruch2019little}.
Our FedGTD protocol computes deviations in the geometric (Riemannian) space
of gradient directions, making it resilient to such coordinated manipulation.
We provide theoretical convergence guarantees under standard smoothness
and bounded-heterogeneity assumptions (Section~\ref{sec:fedgtd}) and
empirically demonstrate superior robustness--utility trade-offs compared
to Krum, trimmed mean, and Bulyan across non-IID data partitions
characteristic of multi-site network deployments.

% === END SEGMENT 1 ===
